\section{おわりに}

本稿では，脆弱性検知に不可欠な検知条件が現在どこから供給されているかを定量的に調査した．
NVDのCPE付与率は2023年公開分までの100\%から2024年以降56.1\%まで低下し，低下分はNVDが明示的に解析対象外とした \texttt{Deferred} に対応する．
しかし付与率は検知精度を予測しない．
組込み部品の脆弱性についてNVDもVulnCheckもFortinet製品のCPEを1件も列挙せず，Fortinetが採番したCVEでも一致は63.3\%にとどまった．
入力データソースのみを変えた4種のデータベースで10環境を検知すると，NVDのみでの未検知率は製品によって1.3\%から68.8\%まで開き，VulnCheckは付与率を上げる一方でWindowsに70.6\%の過剰検知を生んだ．
誤り方も製品で違い，Ciscoでは未検知53.7\%と過剰検知23.2\%が同時に，Appleでは下限を持たないCPEによる過剰検知81.5\%が生じた．
誤りは2023年以前のCVEに限っても減らず，「CVE増加前のNVDなら信頼してよい」という仮説は成り立たない．
併せて，我々自身が置いている抑制規則の代償（Cisco IOS XEでは抑制112件中31件が真陽性）と限界（偽陽性833件のすり抜け）を定量化した．

一次データソースを構造化する手段としては，欠損データの生成・抽出器の生成・HTMLからの直接変換の3パターンを整理し，Fortinet・Apple・Microsoft・OpenSSHに実装した．
今後の課題は，一次データソースから取得すべき製品とNVDやVulnCheckで十分な製品を見分ける方法を，より広い製品で確立することである．取得元を製品ごとに判断できなければ，一次データソースの収集コストは正当化できない．
